\documentclass[10pt]{extarticle}
\usepackage[a4paper,margin=1.5cm]{geometry}
\usepackage[T1]{fontenc}
\usepackage{amsmath,amssymb}
\usepackage{graphicx}
\usepackage[export]{adjustbox}
\usepackage{multicol}
\usepackage{enumitem}
\usepackage{booktabs,array,tabularx}
\setlength{\columnsep}{0.7cm}
\setlength{\parindent}{0pt}
\setlength{\parskip}{3pt}
\pagestyle{empty}
\makeatletter
\newcommand{\qfigdef}[4]{\@namedef{qf@file@#1}{#2}\@namedef{qf@w@#1}{#3}\@namedef{qf@wide@#1}{#4}}
\newcommand{\qf@check}[1]{\@ifundefined{qf@file@#1}{\PackageError{score2}{Unknown figure id #1}{}}{}}
\newcommand{\qf@img}[4][]{\includegraphics[width=#2\dimexpr\@nameuse{qf@w@#4}pt\relax,max width=#3,max height=0.45\textheight,#1]{\@nameuse{qf@file@#4}}}
\newif\ifqf@full
\newif\ifqf@boxed
\newsavebox{\qf@box}
\newcommand{\qf@open}{\global\qf@fulltrue\par\global\qf@boxedtrue
  \setbox\qf@box\vbox\bgroup\hsize\textwidth\linewidth\textwidth}
\newcommand{\qf@emit}[1]{\par\egroup\global\qf@boxedfalse
  \ifdim\dimexpr\ht\qf@box+\dp\qf@box\relax<#1\textheight\noindent\usebox{\qf@box}\par\else\unvbox\qf@box\fi}
\newcommand{\qfullsplit}{\ifqf@boxed\qf@emit{0.85}\fi}
\newcommand{\qf@close}{\ifqf@boxed\qf@emit{0.45}\fi\global\qf@fullfalse}
\newcommand{\qfullbegin}{\end{multicols}\qf@open}
\newcommand{\qfullend}{\qf@close\begin{multicols}{2}}
\newcommand{\qfullstart}{\qf@open}
\newcommand{\qfulldone}{\qf@close}
\newcommand{\qfig}[2][1]{\qf@check{#2}%
  \ifqf@full\qf@blk{#1}{#2}\else
  \ifnum\@nameuse{qf@wide@#2}=1 \expandafter\@firstoftwo\else\expandafter\@secondoftwo\fi
  {\end{multicols}\begin{center}\qf@img{#1}{\textwidth}{#2}\end{center}\begin{multicols}{2}}%
  {\qf@blk{#1}{#2}}\fi}
\newcommand{\qf@blk}[2]{\par\smallskip{\centering\qf@img{#1}{\linewidth}{#2}\par}\smallskip}
\newcommand{\qfigi}[2][1]{\qf@check{#2}\qf@img[valign=c]{#1}{\linewidth}{#2}}
\makeatother
\newcommand{\papersec}[1]{\par\bigskip{\centering\Large\bfseries PAPER-#1\par}\medskip}
\newcommand{\partsec}[2]{\par\medskip\noindent{\large\bfseries PART-#1: #2}\par\smallskip}
\newcommand{\sect}[1]{\par\smallskip\noindent\textbf{#1}\par}
\newcommand{\qn}[1]{\par\smallskip\noindent\textbf{#1.}~}
\newcommand{\paraq}[2]{\par\smallskip\noindent\textbf{Paragraph for Questions #1 and #2}\par}
\newenvironment{qoptions}{\begin{itemize}[leftmargin=1.8em,itemsep=1pt,topsep=2pt,label={}]}{\end{itemize}}
\newcommand{\opt}[2]{\item[(#1)] #2}
% figure registry (generated by score2_tools.py finalize)
\qfigdef{p1_phy_s1_q01_1}{Score-II - Test 41 (19 Mar 2025)-img/p1_phy_s1_q01_1.png}{96.9}{0}
\qfigdef{p1_phy_s1_q04_1}{Score-II - Test 41 (19 Mar 2025)-img/p1_phy_s1_q04_1.png}{59.9}{0}
\qfigdef{p1_phy_s1_q06_1}{Score-II - Test 41 (19 Mar 2025)-img/p1_phy_s1_q06_1.png}{121.6}{0}
\qfigdef{p1_phy_s1_q08_1}{Score-II - Test 41 (19 Mar 2025)-img/p1_phy_s1_q08_1.png}{144.9}{0}
\qfigdef{p1_phy_s1_q09_1}{Score-II - Test 41 (19 Mar 2025)-img/p1_phy_s1_q09_1.png}{123.7}{0}
\qfigdef{p1_phy_s1_q10_1}{Score-II - Test 41 (19 Mar 2025)-img/p1_phy_s1_q10_1.png}{155.3}{0}
\qfigdef{p1_phy_s1_q12_1}{Score-II - Test 41 (19 Mar 2025)-img/p1_phy_s1_q12_1.png}{91.8}{0}
\qfigdef{p1_phy_s1_q13_1}{Score-II - Test 41 (19 Mar 2025)-img/p1_phy_s1_q13_1.png}{106.6}{0}
\qfigdef{p1_phy_s1_q16_1}{Score-II - Test 41 (19 Mar 2025)-img/p1_phy_s1_q16_1.png}{165.6}{0}
\qfigdef{p1_phy_s1_q17_1}{Score-II - Test 41 (19 Mar 2025)-img/p1_phy_s1_q17_1.png}{48.0}{0}
\qfigdef{p1_phy_s1_q18_1}{Score-II - Test 41 (19 Mar 2025)-img/p1_phy_s1_q18_1.png}{98.7}{0}
\qfigdef{p1_phy_s1_q19_1}{Score-II - Test 41 (19 Mar 2025)-img/p1_phy_s1_q19_1.png}{84.6}{0}
\qfigdef{p1_phy_s1_q20_1}{Score-II - Test 41 (19 Mar 2025)-img/p1_phy_s1_q20_1.png}{95.8}{0}
\qfigdef{p1_phy_s2_q01_1}{Score-II - Test 41 (19 Mar 2025)-img/p1_phy_s2_q01_1.png}{129.3}{0}
\qfigdef{p1_phy_s2_q02_1}{Score-II - Test 41 (19 Mar 2025)-img/p1_phy_s2_q02_1.png}{90.0}{0}
\qfigdef{p1_phy_s2_q03_1}{Score-II - Test 41 (19 Mar 2025)-img/p1_phy_s2_q03_1.png}{54.2}{0}
\qfigdef{p1_phy_s2_q05_1}{Score-II - Test 41 (19 Mar 2025)-img/p1_phy_s2_q05_1.png}{137.3}{0}
\qfigdef{p1_chem_s1_q06_1}{Score-II - Test 41 (19 Mar 2025)-img/p1_chem_s1_q06_1.png}{122.6}{0}
\qfigdef{p1_chem_s1_q06_2}{Score-II - Test 41 (19 Mar 2025)-img/p1_chem_s1_q06_2.png}{78.2}{0}
\qfigdef{p1_chem_s1_q06_3}{Score-II - Test 41 (19 Mar 2025)-img/p1_chem_s1_q06_3.png}{85.7}{0}
\qfigdef{p1_chem_s1_q10_1}{Score-II - Test 41 (19 Mar 2025)-img/p1_chem_s1_q10_1.png}{145.1}{0}
\qfigdef{p1_chem_s1_q11_1}{Score-II - Test 41 (19 Mar 2025)-img/p1_chem_s1_q11_1.png}{161.3}{0}
\qfigdef{p1_chem_s1_q11_2}{Score-II - Test 41 (19 Mar 2025)-img/p1_chem_s1_q11_2.png}{41.3}{0}
\qfigdef{p1_chem_s1_q11_3}{Score-II - Test 41 (19 Mar 2025)-img/p1_chem_s1_q11_3.png}{52.4}{0}
\qfigdef{p1_chem_s1_q11_4}{Score-II - Test 41 (19 Mar 2025)-img/p1_chem_s1_q11_4.png}{29.9}{0}
\qfigdef{p1_chem_s1_q11_5}{Score-II - Test 41 (19 Mar 2025)-img/p1_chem_s1_q11_5.png}{70.4}{0}
\qfigdef{p1_chem_s1_q12_1}{Score-II - Test 41 (19 Mar 2025)-img/p1_chem_s1_q12_1.png}{132.9}{0}
\qfigdef{p1_chem_s1_q17_1}{Score-II - Test 41 (19 Mar 2025)-img/p1_chem_s1_q17_1.png}{89.5}{0}
\qfigdef{p1_chem_s2_q05_1}{Score-II - Test 41 (19 Mar 2025)-img/p1_chem_s2_q05_1.png}{95.7}{0}
\qfigdef{p1_math_s1_q08_1}{Score-II - Test 41 (19 Mar 2025)-img/p1_math_s1_q08_1.png}{80.3}{0}
\qfigdef{p1_math_s1_q15_1}{Score-II - Test 41 (19 Mar 2025)-img/p1_math_s1_q15_1.png}{152.1}{0}
\makeatletter\@beginparpenalty=10000\makeatother
\begin{document}
\begin{center}{\Large\bfseries JEE (Main+Advanced) Score-II --- Test 41 (19 Mar 2025)}\end{center}
\begin{multicols}{2}
\raggedcolumns
\partsec{1}{PHYSICS}
\sect{Section-I}

\qn{1} The net force experienced by dipole (placed along the axis of ring at distance $\sqrt2$ R from centre of the ring of Radius R) due to uniformly charged ring as shown in figure is
\qfig{p1_phy_s1_q01_1}
\begin{qoptions}
\opt{A}{Towards left}
\opt{B}{Towards right}
\opt{C}{Upward}
\opt{D}{No force}
\end{qoptions}

\qn{2} \textbf{Assertion (A) :-} The electrostatic force between the plates of a charged isolated capacitor decreases when dielectric just fill the whole space between the plates.\\
\textbf{Reason (R) :-} The electric field between the plates of a charged isolated capacitor decreases when dielectric fills the whole space.
\begin{qoptions}
\opt{A}{Both A and R are true and R is correct explanation of A}
\opt{B}{Both A and R are true but R is not correct explanation of A}
\opt{C}{Assertion A is true but Reason R is false}
\opt{D}{Assertion A is false but Reason R is true}
\end{qoptions}

\qn{3} The electron of hydrogen atom is considered to be revolving around the fixed proton in circular orbit of radius $\dfrac{b^2}{me^2}$ with velocity $\dfrac{e^2}{b}$ where $b=\dfrac{h}{2\pi}$, the electric current associated due to this revolution is
\begin{qoptions}
\opt{A}{$\dfrac{2\pi^2me^5}{h^3}$}
\opt{B}{$\dfrac{4\pi^2me^5}{h^3}$}
\opt{C}{$\dfrac{2\pi\,me^5}{h^3}$}
\opt{D}{$\dfrac{4\pi\,me^5}{h^3}$}
\end{qoptions}

\qn{4} At time $t=0$, the LC circuit shown in the figure has equal amount of energy stored in the capacitor $C=20\ \mu$F and in the inductor, each of 200$\mu$J. Current amplitude is $\dfrac1{\sqrt{10}}$ A. The value of inductor is
\qfig{p1_phy_s1_q04_1}
\begin{qoptions}
\opt{A}{2 mH}
\opt{B}{4 mH}
\opt{C}{8 mH}
\opt{D}{16 mH}
\end{qoptions}

\qn{5} A negatively charged ion (with charge $-q$ and mass $m$) moving with velocity $v_0\hat i$ enters a region of uniform magnetic field of Intensity $B_0\hat k$ at time $t=0$. Here $B_o$ and $v_0$ are positive constants. The angular momentum of ion at time $t=\dfrac{\pi m}{2B_0q}$ about the point where the ion enters the region of field is
\begin{qoptions}
\opt{A}{$\dfrac{m^2v_0^2}{B_0q}\hat k$}
\opt{B}{$\dfrac{m^2v_0^2}{B_0q}\left(-\hat k\right)$}
\opt{C}{zero}
\opt{D}{$\dfrac{m^2v_0^2}{2B_0q}\left(\hat k\right)$}
\end{qoptions}

\qn{6} The instantaneous voltage drop $V_R$ across resistor in circuit is
\qfig{p1_phy_s1_q06_1}
\begin{qoptions}
\opt{A}{$V_R=\dfrac{-V_0}{\sqrt2}\cos t$}
\opt{B}{$V_R=V_0\cos t$}
\opt{C}{$V_R=\dfrac{-V_0}{\sqrt2}\cos\left(t+\dfrac\pi4\right)$}
\opt{D}{$V_R=\dfrac{-V_0}{\sqrt2}\sin\left(t-\dfrac\pi4\right)$}
\end{qoptions}

\qn{7} A parallel plate capacitor (with circular plates of radius $r$ and plate separation $d$) is being charged. Magnitude of electric field in between the plates of capacitor at any time $t$ is $E$. Magnitude of poynting vector at outer edge of capacitor plate just inside the plate is (where $\mu_0$ and $\in_0$ respectively represent magnetic permeability and electrical permittivity of free space)
\begin{qoptions}
\opt{A}{$\dfrac12\in_0rE\left(\dfrac{dE}{dt}\right)$}
\opt{B}{$\dfrac12\mu_0rE\left(\dfrac{dE}{dt}\right)$}
\opt{C}{$\dfrac12\dfrac{\in_0}{\mu_0}rE\left(\dfrac{dE}{dt}\right)$}
\opt{D}{$\dfrac12\dfrac{\mu_0}{\in_0}rE\left(\dfrac{dE}{dt}\right)$}
\end{qoptions}

\qn{8} In the given figure there are two ideal, relaxed springs S$_1$ \& S$_2$ of spring constants $k_1$ \& $k_2$ respectively. An application of external Variable force $F$ causes block to move very slowly towards the wall on smooth horizontal floor by a distance $x$.
\qfig{p1_phy_s1_q08_1}
\par\nobreak
\begin{center}\small
\begin{tabular}{|c|p{2.3cm}|c|p{2.6cm}|}\hline
 & \multicolumn{1}{c|}{\textbf{Column-I}} & & \multicolumn{1}{c|}{\textbf{Column-II}}\\\hline
(I) & Work done by S$_2$ on block & (P) & Zero\\\hline
(II) & Work done by S$_2$ on S$_1$ & (Q) & $-\frac12\frac{k_1k_2}{k_1+k_2}x^2$\\\hline
(III) & Work done by $F$ on block & (R) & $\frac12\frac{k_1k_2}{k_1+k_2}x^2$\\\hline
(IV) & Work done by S$_1$ on wall & (S) & $\frac12\frac{k_1k_2^2}{(k_1+k_2)^2}x^2$\\\hline
\end{tabular}
\end{center}
\begin{qoptions}
\opt{A}{I $-$ Q ; II $-$ S ; III $-$ R ; IV $-$ P}
\opt{B}{I $-$ R ; II $-$ S ; III $-$ Q ; IV $-$ P}
\opt{C}{I $-$ S ; II $-$ R ; III $-$ P ; IV $-$ Q}
\opt{D}{I $-$ R ; II $-$ S ; III $-$ P ; IV $-$ Q}
\end{qoptions}

\qn{9} A light ray AB is incident on solid sphere of glass of refractive index $\sqrt3$. A mirror is placed at some distance which is perpendicular to BO. Where O is centre of sphere and B is incident point of ray, then final emergent ray from glass sphere after reflection from plane mirror is :-
\qfig{p1_phy_s1_q09_1}
\begin{qoptions}
\opt{A}{Horizontal and parallel to BO}
\opt{B}{Vertical and perpendicular to BO}
\opt{C}{Ray will retrace it's path and finally emerge out from point B as BA}
\opt{D}{Light can not emerge out from sphere as T.I.R takes place from sphere to air boundary}
\end{qoptions}

\qn{10} In the following YDSE setup, there is monochromatic light source S of wavelength 5000 \AA\ is placed 1m behind the plane of slits and in front of one of the slits, then central maxima O$'$ is shifted by distance $y$ from central point O as shown in figure then
\qfig{p1_phy_s1_q10_1}
\begin{qoptions}
\opt{A}{$y=1$ mm and I$^{\text{st}}$ maxima at O}
\opt{B}{$y=1$ mm and I$^{\text{st}}$ minima at O}
\opt{C}{$y=0.75$ mm and I$^{\text{st}}$ minima at O}
\opt{D}{$y=0.75$ mm and I$^{\text{st}}$ maxima at O}
\end{qoptions}

\qn{11} If $E_n$ represents the kinetic energy of electron in $n^{\text{th}}$ orbit of hydrogen like atom then select the correct relation from the following.
\begin{qoptions}
\opt{A}{$E_n=\dfrac{E_{n+1}+E_{n-1}}2$}
\opt{B}{$E_n<\dfrac{E_{n+1}+E_{n-1}}2$}
\opt{C}{$E_n>\dfrac{E_{n+1}+E_{n-1}}2$}
\opt{D}{$E_n=\dfrac{E_{n+1}-E_{n-1}}2$}
\end{qoptions}

\qn{12} A graph is plotted between rate of disintegration $\left(\dfrac{-dN}{dt}\right)$ and number of active nuclei at any instant. What is half life of radio active substance? (If time in seconds)
\qfig{p1_phy_s1_q12_1}
\begin{qoptions}
\opt{A}{$\sqrt3$ sec}
\opt{B}{$\dfrac1{\sqrt3}$ sec}
\opt{C}{$\sqrt3\,\ell$n2 sec}
\opt{D}{$\dfrac{\ell\text{n}2}{\sqrt3}$ sec}
\end{qoptions}

\qn{13} The density versus pressure graph of one mole of an ideal monoatomic gas of molar mass $M$ undergoing a process AB is shown in figure. The amount of heat rejected by gas in the process is
\qfig{p1_phy_s1_q13_1}
\begin{qoptions}
\opt{A}{$\dfrac{2P_0M}{\rho_0}\ell$n2}
\opt{B}{$\dfrac{P_0M}{\rho_0}\ell$n2}
\opt{C}{$\dfrac12\dfrac{P_0M}{\rho_0}\ell$n2}
\opt{D}{$4\,\dfrac{P_0M}{\rho_0}\ell$n2}
\end{qoptions}

\qn{14} \textbf{Statement-I :-} Two particles of mass 1 kg and 3 kg move towards each other under their mutual force of attraction. No other force acts on them. When the relative velocity of approach of two particle is 2 m/sec, their centre of mass has velocity of 0.5 m/sec. When the relative velocity of approach increases to 3 m/sec, the velocity of the centre of mass decreases to 0.25 m/sec\\
\textbf{Statement-II :-} The total kinetic energy as seen from ground is $\dfrac12\mu V_{\text{rel}}^2+\dfrac12MV_{\text{c.m}}^2$ and here $\mu$ is reduced mass $\mu=\dfrac{m_1m_2}{m_1+m_2}$ and $M$ is total mass $M=(m_1+m_2)$ and in absence of external force, total mechanical energy remains conserved
\begin{qoptions}
\opt{A}{Statement-I is true, Statement-II is true}
\opt{B}{Statement-I is false, Statement-II is true}
\opt{C}{Statement-I is false, Statement-II is false}
\opt{D}{Statement-I is true, Statement-II is false}
\end{qoptions}

\qn{15} A particle of mass $m$ is moving in a horizontal circle of Radius $R$ under a centripetal force given by $\left(F_{\text{cen}}=\dfrac{-2}{R^2}\right)$ then. (If potential energy at infinity is zero)
\begin{qoptions}
\opt{A}{Kinetic energy is $\dfrac1{6R}$}
\opt{B}{Total energy of particle is $\dfrac{-1}R$}
\opt{C}{Kinetic energy of particle is $\dfrac2R$}
\opt{D}{Potential energy is $\dfrac1R$}
\end{qoptions}

\qn{16} A particle of mass $m$ strikes a smooth disc of radius $R$ and of same mass $m$ with velocity $v$. Initially the disc is at rest. The coefficient of restitution $e=0.5$. Find the speed of particle just after the impact. Assuming the disc is placed on smooth horizontal table as shown by top view of table
\qfig{p1_phy_s1_q16_1}
\begin{qoptions}
\opt{A}{$\dfrac{3v}4\sqrt{19}$}
\opt{B}{$\dfrac v4\sqrt{19}$}
\opt{C}{$\dfrac v2\sqrt{19}$}
\opt{D}{$\dfrac v8\sqrt{19}$}
\end{qoptions}

\qn{17} The time period of following oscillation system is :-\\
(All pulley and spring are ideal)
\qfig{p1_phy_s1_q17_1}
\begin{qoptions}
\opt{A}{$\dfrac\pi2\sqrt{\dfrac mk}$}
\opt{B}{$2\pi\sqrt{\dfrac m{32k}}$}
\opt{C}{$2\pi\sqrt{\dfrac{16m}k}$}
\opt{D}{$4\pi\sqrt{\dfrac mk}$}
\end{qoptions}

\qn{18} A ball is projected on a smooth horizontal surface towards a smooth circular track of radius $R$ with speed $\sqrt{5gR}$. The ball after leaving the semi-circular track lands on the smooth ground where coefficient of restitution is $\dfrac12$. The range between first collision and second collision is
\qfig{p1_phy_s1_q18_1}
\begin{qoptions}
\opt{A}{$R$}
\opt{B}{$\dfrac{3R}2$}
\opt{C}{$\dfrac{2R}3$}
\opt{D}{$2R$}
\end{qoptions}

\qn{19} A rigid ring A and a rigid thin disk B both made of same material, when gently placed on water, just manage to float due to surface tension as shown in the figure. Both the ring and the disk have same radius. What can you conclude about their masses (thickness may be different)?
\qfig{p1_phy_s1_q19_1}
\begin{qoptions}
\opt{A}{Both have the same mass.}
\opt{B}{Mass of the ring is half of that of the disk.}
\opt{C}{Mass of the ring is double to that of the disk.}
\opt{D}{More information is needed to decide.}
\end{qoptions}

\qn{20} Three tubes A, B, C are connected to a horizontal pipe in which liquid is flowing. The radii of the pipes at the joints of A, B and C are 2 cm, 1 cm and 2 cm respectively. The height of the liquid
\qfig{p1_phy_s1_q20_1}
\begin{qoptions}
\opt{A}{in A is lesser than C}
\opt{B}{in A and C is equal}
\opt{C}{is same in all the three}
\opt{D}{in A and B is same}
\end{qoptions}

\sect{Section-II}

\qn{1} For the situation shown in figure liquid of coefficient of viscosity $\eta$ flown over the surface of fixed plate in $x$-$z$ plane. The velocity of liquid as a function of $y$ (perpendicular distance from plate to free surface) is given as $v=\alpha\left[\dfrac yh-\left(\dfrac yh\right)^2\right]$. If magnitude of shear stress that liquid applies at base plate is $\dfrac{N\alpha\eta}h$ then $N$ is
\qfig{p1_phy_s2_q01_1}

\qn{2} A solid sphere is released from a point at a height $h=7$m on a parabolic surface as shown in figure. Left part of the surface ($x<0$) is sufficiently rough to avoid slipping and right part of the surface ($x>0$) is perfectly smooth. The maximum height in meter upto which the sphere will reach on right part, is
\qfig{p1_phy_s2_q02_1}

\qn{3} A uniform rod of mass $M$, length $\ell$ and area of cross section $A$ is hanging vertically from ceiling as shown. If Young's modulus of elasticity is $Y$, the elastic potential energy stored in upper half of the rod, is $\dfrac{14M^2g^2\,\ell}{nAY}$, then value of $n$ is \underline{\hspace{2em}}.
\qfig{p1_phy_s2_q03_1}

\qn{4} A reversible heat engine has an efficiency of 33.3\%, removing heat from a hot reservoir and rejecting heat to a cold reservoir at 0$^\circ$C. If the engine now operates in reverse, how long would it take of freeze 1.0 kg of water at 0$^\circ$C, if it operates on a power of 186 W?

\qn{5} The value of power dissipated across the zener diode ($V_z=15$ V) connected in the circuit as shown in the figure is $x\times10^{-1}$ watt.
\qfig{p1_phy_s2_q05_1}
The value of $x$, to the nearest integer, is \underline{\hspace{2em}}.

\partsec{2}{CHEMISTRY}
\sect{Section-I}

\qn{1} Which of the following statement is correct regarding the line spectrum of Hydrogen \& Hydrogen like species ?
\begin{qoptions}
\opt{A}{All the lines of Balmer series for H$-$atom belongs to the visible region.}
\opt{B}{The first line of Balmer series for H$-$atom will be red line.}
\opt{C}{Balmer series for He$^+$ ion belongs to the visible region.}
\opt{D}{Second line of Balmer series of H$-$atom is red line.}
\end{qoptions}

\qn{2} Which of the following Processes are exothermic ?\\
(1) O$^{-}$(g) $\to$ O$^{+}$(g)\\
(2) Be$^{-}$(g) $\to$ Be(g)\\
(3) P(g) $\to$ P$^{3-}$ (g)\\
(4) H$^{+}$(g) $\to$ H$^{-}$(g)
\begin{qoptions}
\opt{A}{1,2,4}
\opt{B}{2,4}
\opt{C}{2,3}
\opt{D}{1,3}
\end{qoptions}

\qn{3} \textbf{Assertion :-} Solubility of Pb(OH)$_2$ (weak base) decreases on addition of NaOH.\\
\textbf{Reason :-} NaOH increases OH$^{-}$ concentration forcing the equilibrium to shift in the backward direction.
\begin{qoptions}
\opt{A}{Assertion and Reason both are correct and reason is correct explanation of assertion.}
\opt{B}{Assertion and Reason both are correct and reason is incorrect explanation of assertion.}
\opt{C}{Assertion and Reason both are incorrect.}
\opt{D}{Assertion is correct and Reason is incorrect.}
\end{qoptions}

\qn{4} Which of the following statement is correct?
\begin{qoptions}
\opt{A}{As per MOT, double bond in C$_2$ consist of sigma bond and pi bond.}
\opt{B}{Bond energy : N$_2$ $>$ C$_2$ $>$ O$_2$ $>$ B$_2$ $>$ F$_2$}
\opt{C}{Bond length : N$_2$ $<$ C$_2$ $<$ O$_2$ $<$ B$_2$ $<$ F$_2$}
\opt{D}{Diamagnetic Li$_2$ molecules are known to exist in solid phase.}
\end{qoptions}

\qn{5} The aqueous solution of which of the following will have highest freezing point?
\begin{qoptions}
\opt{A}{Glucose (0.1 m)}
\opt{B}{NaCl (0.1 m)}
\opt{C}{Na$_3$PO$_4$ (0.05 m)}
\opt{D}{CH$_3$COOH (0.2 m, degree of dissociation = 0.1)}
\end{qoptions}

\qn{6}\par\nobreak
\qfig{p1_chem_s1_q06_1}
The structure of E is :-
\begin{qoptions}
\opt{A}{\qfigi{p1_chem_s1_q06_2}}
\opt{B}{CH$_3$CH$_2$OH}
\opt{C}{\qfigi{p1_chem_s1_q06_3}}
\opt{D}{CH$_3$CH$_2$NH$_2$}
\end{qoptions}

\qn{7} Two system PCl$_5$(g) $\rightleftharpoons$ PCl$_3$(g) + Cl$_2$(g) and COCl$_2$(g) $\rightleftharpoons$ CO(g) + Cl$_2$(g) are simultaneously in equilibrium in a vessel at constant volume. If some CO is introduced in the vessel, then at new equilibrium, the moles of :
\begin{qoptions}
\opt{A}{PCl$_5$ increase}
\opt{B}{PCl$_5$ remain unchanged}
\opt{C}{PCl$_5$ decrease}
\opt{D}{Cl$_2$ increase}
\end{qoptions}

\qn{8} During estimation of nitrogen present in an organic compound by kjeldahl's method, the ammonia evolved from `$m$' g of the compound in kjeldahl's estimation of nitrogen, was absorbed in `$V$'ml of `$M$' molar H$_2$SO$_4$ solution. The excess of acid required `$V_1$'ml of NaOH of molarity `$M$'. The \% of `N' in the organic sample is :-
\begin{qoptions}
\opt{A}{$\dfrac{28\times M\times2\left(V-\frac{V_1}2\right)}{1000}\times\dfrac{100}m$}
\opt{B}{$\dfrac{14\times M\times\left(V-\frac{V_1}2\right)}{1000}\times\dfrac{100}m$}
\opt{C}{$\dfrac{14\times M\times2\,(2V-V_1)}{1000}\times\dfrac{100}m$}
\opt{D}{$\dfrac{28\times M\times\left(V-\frac{V_1}2\right)}{1000}\times\dfrac{100}m$}
\end{qoptions}

\qn{9} Match the column and choose the correct option :
\par\nobreak
\begin{center}
\begin{tabular}{|c|c|c|c|}\hline
\multicolumn{2}{|c|}{\shortstack{\textbf{Column-I}\\\textbf{(multiple)}}} & \multicolumn{2}{c|}{\shortstack{\textbf{Column-II}\\\textbf{(Prefix)}}}\\\hline
(A) & $10^9$ & (1) & atto\\\hline
(B) & $10^{21}$ & (2) & zeta\\\hline
(C) & $10^{-18}$ & (3) & exa\\\hline
(D) & $10^{18}$ & (4) & giga\\\hline
\end{tabular}
\end{center}
\begin{qoptions}
\opt{A}{A $-$ 2 ; B $-$ 1 ; C $-$ 3 ; D $-$ 4}
\opt{B}{A $-$ 4 ; B $-$ 2 ; C $-$ 1 ; D $-$ 3}
\opt{C}{A $-$ 1 ; B $-$ 2 ; C $-$ 3 ; D $-$ 4}
\opt{D}{A $-$ 3 ; B $-$ 2 ; C $-$ 4 ; D $-$ 1}
\end{qoptions}

\qn{10}\par\nobreak
\qfig{p1_chem_s1_q10_1}
The galvanometer shows no deflection when $R_1=100$ ohm. Find the molar conductivity of 0.05 M KCl solution, if the cell constant is 1.25 cm$^{-1}$.
\begin{qoptions}
\opt{A}{200 S cm$^2$ mol$^{-1}$}
\opt{B}{400 S cm$^2$ mol$^{-1}$}
\opt{C}{1000 S cm$^2$ mol$^{-1}$}
\opt{D}{125 S cm$^2$ mol$^{-1}$}
\end{qoptions}

\qn{11}\par\nobreak
\qfig{p1_chem_s1_q11_1}
What is the correct structure of D is :-
\begin{qoptions}
\opt{A}{\qfigi{p1_chem_s1_q11_2}}
\opt{B}{\qfigi{p1_chem_s1_q11_3}}
\opt{C}{\qfigi{p1_chem_s1_q11_4}}
\opt{D}{\qfigi{p1_chem_s1_q11_5}}
\end{qoptions}

\qn{12} Give the correct order of Acidic strength.
\qfig{p1_chem_s1_q12_1}
\begin{qoptions}
\opt{A}{(1) $>$ (2) $>$ (3) $>$ (4)}
\opt{B}{(3) $>$ (1) $>$ (2) $>$ (4)}
\opt{C}{(2) $>$ (1) $>$ (4) $>$ (3)}
\opt{D}{(2) $>$ (1) $>$ (3) $>$ (4)}
\end{qoptions}

\qn{13} Consider the complex : $[\mathrm{Ni(NH_3)_{6-2n}(en)_n}]^{2+}$ The maximum number of sterioisomers is possible for -
\begin{qoptions}
\opt{A}{$n=0$}
\opt{B}{$n=1$}
\opt{C}{$n=2$}
\opt{D}{$n=3$}
\end{qoptions}

\qn{14} For the complex $[\mathrm{Co(C_2O_4)_3}]^{3-}$, which information is ``Incorrect'' ?
\begin{qoptions}
\opt{A}{$\mathrm{C_2O_4^{2-}}$ is monodentate ligand.}
\opt{B}{Inner orbital complex.}
\opt{C}{Oxidation state of Co = +3.}
\opt{D}{Complex has octahedral geometry.}
\end{qoptions}

\qn{15} \textbf{Statement-1 :} Benzaldehyde gives positive test with both Tollens' reagent and Fehling's reagent.\\
\textbf{Statement-2 :} Ph $-$ CH$_2$ $-$ CHO gives positive test with Fehling's reagent.
\begin{qoptions}
\opt{A}{Statement-1 is true, statement-2 is false.}
\opt{B}{Statement-1 is false, statement-2 is true.}
\opt{C}{Both Statement-1 and statement-2 are true.}
\opt{D}{Both Statement-1 and statement-2 are false.}
\end{qoptions}

\qn{16} A compound `X' having molecular formula C$_4$H$_{11}$N give reaction with C$_6$H$_5$SO$_2$Cl and give compounds `Y' and HCl, But does not give reaction with chloroform in presence of KOH. Compound `X' also undergo acylation with ethanoyl chloride to yield `Z'.\\
Choose incorrect statement.
\begin{qoptions}
\opt{A}{Compound `Y' is insoluble in NaOH.}
\opt{B}{`Z' is an Amide.}
\opt{C}{`X' reacts with nitrous acid to give N$_2$ gas.}
\opt{D}{Compound `X' does not give mustard oil test.}
\end{qoptions}

\qn{17} What is the common name of the following compound ?
\qfig{p1_chem_s1_q17_1}
\begin{qoptions}
\opt{A}{Tricarboxylic acid}
\opt{B}{Carballylic acid}
\opt{C}{Propane tricarboxylic acid}
\opt{D}{Hexane trioic acid}
\end{qoptions}

\qn{18} What will be the expression of rate constant ($k$) of reaction in terms of pressure for following reaction, when the initial pressure is $P_0$ and, at time `$t$', total pressure is $P_T$ (considering reaction is following first order kinetics) :\\
\centerline{2A(g) $\longrightarrow$ B(g)}
\begin{qoptions}
\opt{A}{$k=\dfrac1t\,\ell\text{n}\dfrac{P_0}{2P_0-P_T}$}
\opt{B}{$k=\dfrac1{2t}\,\ell\text{n}\dfrac{P_0}{P_0-P_T}$}
\opt{C}{$k=\dfrac1{2t}\,\ell\text{n}\dfrac{P_0}{2P_T-P_0}$}
\opt{D}{$k=\dfrac{2.303}t\,\ell\text{og}\dfrac{P_0}{P_T}$}
\end{qoptions}

\qn{19} Calculate enthalpy of reaction for following reaction -\\
C$_3$H$_8$(g) + H$_2$(g) $\longrightarrow$ CH$_4$(g) + C$_2$H$_6$(g)\\
Given that -\\
Enthalpy of combustions :\\
$\Delta_c$H$^\text{o}$ (Graphite, s) $=-393$ kJ mol$^{-1}$\\
$\Delta_c$H$^\text{o}$ (CH$_4$, g) $=-890$ kJmol$^{-1}$\\
$\Delta_c$H$^\text{o}$ (C$_2$H$_6$, g) $=-1560$ kJmol$^{-1}$\\
and enthalpy of formation :\\
$\Delta_f$H$^\text{o}$ (Propane, g) $=-104$ kJmol$^{-1}$\\
$\Delta_f$H$^\text{o}$ (H$_2$O, $\ell$) $=-285$ kJmol$^{-1}$
\begin{qoptions}
\opt{A}{$-50$ kJmol$^{-1}$}
\opt{B}{$+2061$ kJmol$^{-1}$}
\opt{C}{$-2061$ kJmol$^{-1}$}
\opt{D}{$+50$ kJmol$^{-1}$}
\end{qoptions}

\qn{20} What is the f-subshell electronic configuration of Am$^{3+}$ (Z = 95)?
\begin{qoptions}
\opt{A}{5f$^4$}
\opt{B}{5f$^5$}
\opt{C}{5f$^6$}
\opt{D}{5f$^7$}
\end{qoptions}

\sect{Section-II}

\qn{1} According to VSEPR theory, how many single super pairs are present in C$_2$(CN)$_4$

\qn{2} Iron (III) chloride (aq.) is treated with potassium ferrocyanide to give a blue ppt `A'. Which on further treatment with excess of potassium ferrocyanide give a colloidal solution `B'.\\
Find the sum of number of Iron atoms present in `A' and `B'.

\qn{3} Find the molecular mass of Tripeptide formed using Amino acids $-$ Glycine, Alanine, Leucine.

\qn{4} How many of the following compound will give turbidity in 5 to 10 minutes with Lucas Reagent ?\\
Methyl Alcohol, n-Propyl alcohol, Isopropyl alcohol, n-Butyl alcohol, sec-Butyl alcohol, Isobutyl alcohol, tert-Butyl alcohol.

\qn{5}\par\nobreak
\qfig{p1_chem_s2_q05_1}
Two flasks `1' \& `2' are filled with H$_2$ gas as shown above connected by a valve of negligible volume. When the valve is opened, what will be the change in entropy, (in cal K$^{-1}$) due to mixing.\\
Given : ln 2 = 0.7 R = 2 cal mol$^{-1}$ K$^{-1}$

\partsec{3}{MATHEMATICS}
\sect{Section-I}

\qn{1} Sum of series $2+\dfrac57+\dfrac8{49}+\dfrac{11}{343}+\dots.\infty$ is
\begin{qoptions}
\opt{A}{$\dfrac{35}{12}$}
\opt{B}{$\dfrac{49}{18}$}
\opt{C}{$\dfrac{35}{16}$}
\opt{D}{$\dfrac7{12}$}
\end{qoptions}

\qn{2} Given a right circular cone has a volume $V_1$ and the largest right circular cylinder that can be inscribed in the cone has a volume $V_2$. Then $V_1:V_2$ is $\dfrac ab$ (where $a,b$ are coprime). Find the value of $a^2-b^2$
\begin{qoptions}
\opt{A}{25}
\opt{B}{65}
\opt{C}{64}
\opt{D}{20}
\end{qoptions}

\qn{3} \textbf{Assertion (A) :} Line $7x+5y+1=0$ cannot intersect the parabola $y=x^2$ in a rational point (whose abscissa \& ordinate both are rational)\\
\textbf{Reason (R) :} For odd integer $a,b$ and $c$ equation $ax^2+bx+c=0$ cannot have rational roots.
\begin{qoptions}
\opt{A}{Both \textbf{A} and \textbf{R} are correct and \textbf{R} is the correct explanation of \textbf{A}}
\opt{B}{Both \textbf{A} and \textbf{R} are correct but \textbf{R} is \textbf{NOT} the correct explanation of \textbf{A}}
\opt{C}{\textbf{A} is correct but \textbf{R} is not correct}
\opt{D}{\textbf{A} is not correct but \textbf{R} is correct}
\end{qoptions}

\qn{4} Let $f:(1,\infty)\to(0,\infty)$ be a continuous decreasing function with $\displaystyle\lim_{x\to\infty}\dfrac{f(x)}{f(8x)}=1$ Then $\displaystyle\lim_{x\to\infty}\dfrac{f(5x)}{f(8x)}$ is equal to
\begin{qoptions}
\opt{A}{$\dfrac48$}
\opt{B}{$\dfrac46$}
\opt{C}{$\dfrac68$}
\opt{D}{1}
\end{qoptions}

\qn{5} Let a ray of light be travelling along the, line : $4x+3y-2k=0$ and strike, a parabolic mirror at $(1,p)$. After reflection it moves, along the : line $3x-4y+k=0$. If the directrix of the parabola is the line $3x-4y+15=0$ then the length of latus-rectum of the parabola is :
\begin{qoptions}
\opt{A}{4}
\opt{B}{6}
\opt{C}{8}
\opt{D}{13}
\end{qoptions}

\qn{6} If AB is a chord of the circle $x^2+y^2=25$. The tangents at A and B intersect at C. If (2,3) is the mid point of AB, then the area of the quadrilateral OACB is, where O is the origin.
\begin{qoptions}
\opt{A}{$\sqrt{\dfrac3{13}}$ sq. unit}
\opt{B}{$50\sqrt{\dfrac3{13}}$ sq. unit}
\opt{C}{100 sq. unit}
\opt{D}{$100\sqrt3$ sq. unit}
\end{qoptions}

\qn{7} A company has 10 employees. The company has decided to form a team including at least three employees and also excluding at least three employees. Then the number of ways of forming the team is
\begin{qoptions}
\opt{A}{950}
\opt{B}{912}
\opt{C}{1000}
\opt{D}{953}
\end{qoptions}

\qn{8} A ray of light passing through the point A(1, 2) is reflected at a point B on the $x-$axis and then passes through (5, 3). Then the equation of AB is :
\qfig{p1_math_s1_q08_1}
\begin{qoptions}
\opt{A}{$5x+4y=13$}
\opt{B}{$5x-4y=-3$}
\opt{C}{$4x+5y=14$}
\opt{D}{$4x-5y=-6$}
\end{qoptions}

\qn{9}
\par\nobreak
\begin{center}\small
\begin{tabular}{|c|p{3.7cm}|c|c|}\hline
 & \multicolumn{1}{c|}{\textbf{List-I}} & & \textbf{List-II}\\\hline
(A) & If $\alpha,\beta$ ($\alpha<\beta$) are the two roots of the equation $\dfrac{1-8(\log_{10}x)^2}{\log_{10}x-2(\log_{10}x)^2}=1$ the the value of $\dfrac{(\beta^3+1)}{\alpha^2\beta^3}$ is equal to & (I) & 5\\\hline
(B) & Number of solution(s) of equation $\log_2(x^2+3)=\dfrac12\log_{1/3}\left(x+\dfrac1x\right)$, $x>0$ is & (II) & 11\\\hline
(C) & Integers satisfying the equation $|x|+\left|\dfrac{4-x^2}{x}\right|=\left|\dfrac4x\right|$ is/are & (III) & 0\\\hline
(D) & If $\log_c2\log_b125=\log_{10}8\log_c10$ where $c>0$, $c\neq1$, $b>0$, $b\neq1$, then $b$ is & (IV) & 1\\\hline
\end{tabular}
\end{center}
\begin{qoptions}
\opt{A}{A $\to$ I, B $\to$ III, C $\to$ IV, D $\to$ II}
\opt{B}{A $\to$ II, B $\to$ III, C $\to$ IV, D $\to$ I}
\opt{C}{A $\to$ IV, B $\to$ II, C $\to$ IV, D $\to$ I}
\opt{D}{A $\to$ IV, B $\to$ III, C $\to$ II, D $\to$ I}
\end{qoptions}

\qn{10} The A.M. of $n$ observations is $M$. If the sum of $(n-4)$ observation is $a$, then the mean of remaining 4 observations is ($a\neq nM$) :
\begin{qoptions}
\opt{A}{$\dfrac{nM-a}4$}
\opt{B}{$\dfrac{nM+a}4$}
\opt{C}{$\dfrac{nM-a}2$}
\opt{D}{$nM+a$}
\end{qoptions}

\qn{11} If $\displaystyle\int\dfrac{(\sqrt x)^5}{(\sqrt x)^7+x^6}dx=\dfrac1k\ln\left(\dfrac{x^k}{1+x^k}\right)+c$ then $k$ is equal to :\\
(where $c$ is constant of integration)
\begin{qoptions}
\opt{A}{$\dfrac15$}
\opt{B}{$\dfrac25$}
\opt{C}{5}
\opt{D}{$\dfrac52$}
\end{qoptions}

\qn{12} Number of roots of the function $f(x)=\dfrac1{(x+1)^3}-3x+\sin x$ is :
\begin{qoptions}
\opt{A}{0}
\opt{B}{1}
\opt{C}{2}
\opt{D}{More than 2}
\end{qoptions}

\qn{13} The value of the expression,\\
$14\tan\big(\tan^{-1}\tfrac13+\tan^{-1}\tfrac17+\tan^{-1}\tfrac1{13}$\\$+\tan^{-1}\tfrac1{21}+\tan^{-1}\tfrac1{31}\big)$ is an integer which is equal to :
\begin{qoptions}
\opt{A}{2}
\opt{B}{5}
\opt{C}{7}
\opt{D}{10}
\end{qoptions}

\qn{14} $\displaystyle\lim_{x\to0}\dfrac{\left(e^{2x^2}-1-2x^2\right)(\cos x-1)}{(\sin3x-\ln(1+3x))\,x^4}$ :
\begin{qoptions}
\opt{A}{$\dfrac29$}
\opt{B}{$\dfrac19$}
\opt{C}{$-\dfrac19$}
\opt{D}{$-\dfrac29$}
\end{qoptions}

\qn{15} Let $f(x)$ be a continous function in $(0,7)$ and graph of $f'(x)$ is given below
\qfig{p1_math_s1_q15_1}
$S_1$ : $f(x)$ has three points of extrema\\
$S_2$ : $f(x)$ has 4 points of inflection
\begin{qoptions}
\opt{A}{Both statement-I and statement-II are correct.}
\opt{B}{Both statement-I and statement-II are incorrect.}
\opt{C}{Statement-I is incorrect and statement-II is correct.}
\opt{D}{Statement-I is correct and statement-II is incorrect.}
\end{qoptions}

\qn{16} If the ellipse $\dfrac{x^2}{16}+\dfrac{y^2}7=1$ and the hyperbola $\dfrac{x^2}{144}-\dfrac{y^2}\alpha=\dfrac1{25}$ are orthogonal, then find the value of $4e_1$. (where $e_1$ is the eccentricity of the hyperbola)
\begin{qoptions}
\opt{A}{10}
\opt{B}{5}
\opt{C}{8}
\opt{D}{6}
\end{qoptions}

\qn{17} If $f(x)=2^x+x^3+1$ and $g(x)=f^{-1}(x)$ then value of $\dfrac{g'(2)}{g''(2)}$ lies in
\begin{qoptions}
\opt{A}{Set of irrational numbers}
\opt{B}{Set of prime numbers}
\opt{C}{$[-10,10]$}
\opt{D}{None of these}
\end{qoptions}

\qn{18} $\displaystyle\lim_{t\to0}\dfrac{12^t-3^t-4^t+1}{\sqrt{2\cos t+7}-3}$ is equal to
\begin{qoptions}
\opt{A}{$\log_e12$}
\opt{B}{$\log_e3.\log_e4$}
\opt{C}{$\left(\log_e3.\log_e4\right)^{-1}$}
\opt{D}{$-6\log_e3.\log_e4$}
\end{qoptions}

\qn{19} Let $\vec a,\vec b$ be two non collinear unit vectors and $\vec x=\vec a-\left(\vec a.\vec b\right)\vec b$ and $\vec y=\vec a\times\vec b$. Now consider the following statements :\\
\textbf{Statement-I :} $|\vec x|=|\vec y|$\\
\textbf{Statement-II :} $|\vec y|=|\vec x|+\left|\vec x.\vec b\right|$
\begin{qoptions}
\opt{A}{Statement-I is true but Statement-II is false}
\opt{B}{Both Statement-I and Statement-II are false}
\opt{C}{Both Statement-I and Statement-II are true}
\opt{D}{Statement-I is false but Statement-II is true}
\end{qoptions}

\qn{20} If $I_1=\displaystyle\int_0^{\pi/2}\dfrac{\sqrt{\cot x}}{\sqrt{\cot x}+\sqrt{\tan x}}\,dx$ \& $I_2=\displaystyle\int_{-\pi}^{\pi}\dfrac{\cos^2x}{1+\pi^x}\,dx$, then -
\begin{qoptions}
\opt{A}{$\dfrac1{I_1}-\dfrac1{I_2}=\dfrac3\pi$}
\opt{B}{$\dfrac2{I_1}-\dfrac1{I_2}=\dfrac5\pi$}
\opt{C}{$\dfrac2{I_1}+\dfrac1{I_2}=\dfrac8\pi$}
\opt{D}{$\dfrac1{I_1}+\dfrac1{I_2}=\dfrac6\pi$}
\end{qoptions}

\sect{Section-II}

\qn{1} The probability of getting either all heads or all tails for exactly the second time in the 3$^{\text{rd}}$ trial, if in each trial three coins are tossed, is $p$ then the value of $(64p)$ equals to -

\qn{2} The value of $\displaystyle\int_{-3}^{-1}(x+1)^4(x+2)^3(x+3)^4dx$ is

\qn{3} Number of roots of the equation $\cos^2x+\dfrac{\left(\sqrt3+1\right)}{2}\sin x-\dfrac{\sqrt3}4-1=0$ which lie in the interval $[-\pi,\pi]$ is

\qn{4} The smallest possible natural number `$p$' for which the equation $x^2-px+1258=0$ has integral roots ?

\qn{5} If $x,y,z$ be three distinct real numbers, and $x\neq1$, $y\neq1$, $z\neq1$. If the vectors $\vec p=x\hat i+\hat j+\hat k$, $\vec q=\hat i+y\hat j+\hat k$, $\vec r=\hat i+\hat j+z\hat k$ are such that $(\vec p\times\vec q)\cdot\vec r=0$, then $\dfrac1{1-x}+\dfrac1{1-y}+\dfrac1{1-z}$ is equal to

\end{multicols}
\end{document}
